\subsection{Continuity of the Fourier transform}

PROPOSITION 8. --- Let E be a locally convex space, \(\mu\) a promeasure on E, and \(\Phi\) the Fourier transform of \(\mu\) . One has the inequalities

\[
\left| \Phi (x ^ {\prime}) \right| \leqslant \Phi (0)\tag{53}
\]

\[
\left| \Phi (x ^ {\prime}) - \Phi (y ^ {\prime}) \right| ^ {2} \leqslant 2 \Phi (0) \big (\Phi (0) - \mathcal {R} \Phi (x ^ {\prime} - y ^ {\prime}) \big)\tag{54}
\]

for \(x', y'\) in \(\mathbf{E}'\) .

Formula (5) of No.~3 permits reducing to the case that E is finite-dimensional and \(\mu\) is a measure. Then

\[
\left| \Phi (x ^ {\prime}) \right| = \left| \int_ {\mathrm{E}} e ^ {i \langle x, x ^ {\prime} \rangle} d \mu (x) \right| \leqslant \int_ {\mathrm{E}} | e ^ {i \langle x, x ^ {\prime} \rangle} | d \mu (x) = \int_ {\mathrm{E}} d \mu (x) = \Phi (0),
\]

whence (53). Moreover, if \(a\) and \(b\) are real numbers, then

\[
\left| e ^ {i a} - e ^ {i b} \right| ^ {2} = \left| e ^ {i b} \right| ^ {2} \left| e ^ {i (a - b)} - 1 \right| ^ {2} = \left(e ^ {i (a - b)} - 1\right) \left(e ^ {- i (a - b)} - 1\right) = 2 - 2 \cos (a - b);
\]

by the Cauchy--Schwarz inequality, we then have

\[
\begin{array}{r l} \big | \Phi (x ^ {\prime}) - \Phi (y ^ {\prime}) \big | ^ {2} & = \Big | \int_ {\mathrm{E}} \big (e ^ {i \langle x, x ^ {\prime} \rangle} - e ^ {i \langle x, y ^ {\prime} \rangle} \big) d \mu (x) \Big | ^ {2} \\ & \leqslant \int_ {\mathrm{E}} \big | e ^ {i \langle x, x ^ {\prime} \rangle} - e ^ {i \langle x, y ^ {\prime} \rangle} \big | ^ {2} d \mu (x) \int_ {\mathrm{E}} 1 ^ {2} d \mu (x) \\ & = \int_ {\mathrm{E}} (2 - 2 \cos \langle x, x ^ {\prime} - y ^ {\prime} \rangle) d \mu (x) \cdot \Phi (0) \\ & = 2 \Phi (0) \big (\Phi (0) - \mathcal {R} \Phi (x ^ {\prime} - y ^ {\prime}) \big), \end{array}
\]

whence (54).

COROLLARY. --- Equip E' with a topology compatible with its vector space structure. For \(\Phi\) to be continuous, it is necessary and sufficient that its real part \(\mathcal{R}\Phi\) be continuous at the origin, in which case \(\Phi\) is uniformly continuous.

This follows from the inequality (54).

Let F be a locally convex space. We equip the dual \(F'\) of F with a topology compatible with the duality between F and \(F'\) , and we identify F with the dual of \(F'\) . Consequently, the Fourier transform of a bounded measure \(\mu\) on \(F'\) is the function \(F\mu\) on F defined by

\[
(\mathcal {F} \mu) (x) = \int_ {\mathrm{F} ^ {\prime}} e ^ {i \langle x, x ^ {\prime} \rangle} d \mu (x ^ {\prime}).
\]

PROPOSITION 9. --- If F is barreled, then the Fourier transform of every bounded measure on \(F'\) is a uniformly continuous function on F.

Let \(\mu\) be a bounded measure on \(\mathbf{F}'\) and \(\Phi\) its Fourier transform. Let \(\varepsilon > 0\) . There exists a compact subset \(\mathbf{K}\) of \(\mathbf{F}'\) such that \(\mu(\mathbf{F}' - \mathbf{K}) \leqslant \varepsilon\) . Now, \(\mathbf{K}\) is compact for the weak topology \(\sigma(\mathbf{F}', \mathbf{F})\) , hence is equicontinuous because \(\mathbf{F}\) is barreled (TVS, III, §4, No.~2, Th. 1). Therefore there exists a symmetric neighborhood \(\mathbf{U}\) of \(0\) in \(\mathbf{F}\) whose polar \(\mathbf{U}^{\circ}\) contains \(\mathbf{K}\) . Let \(x\) be in \(\varepsilon \mathrm{U}\) ; then

\[
\Phi (0) - \mathcal {R} \Phi (x) = \int_ {\mathrm{F} ^ {\prime}} \left(1 - \cos \langle x, x ^ {\prime} \rangle\right) d \mu (x ^ {\prime}).
\]

Now, \(0 \leqslant 1 - \cos \langle x, x' \rangle \leqslant 2\) for every \(x' \in \mathbf{F}' - \mathbf{K}\) , and

\[
1 - \cos \langle x, x ^ {\prime} \rangle \leqslant \frac {1}{2} \langle x, x ^ {\prime} \rangle^ {2} \leqslant \frac {\varepsilon^ {2}}{2}
\]

for \(x^{\prime}\in \mathbf{K}\subset \mathbf{U}^{\circ}\) ; it follows that

\[
0 \leqslant \Phi (0) - \mathcal {R} \Phi (x) \leqslant 2 \mu (\mathrm{F} ^ {\prime} - \mathrm{K}) + \frac {\varepsilon^ {2}}{2} \mu (\mathrm{K}) \leqslant 2 \varepsilon + \frac {\varepsilon^ {2}}{2} \mu (\mathrm{F} ^ {\prime}).
\]

The second member of this inequality tends to 0 with \(\varepsilon\) ; thus \(\mathcal{R}\Phi\) is continuous at 0 and the proposition follows from the Cor. of Prop. 8.

